\section{Conclusion}
DisasterClaw connects disaster records, an interactive geographic world, a task benchmark, and traceable closed-loop inspection. Paired xBD images and building annotations are anchored to map coordinates shared by the vehicle, observation footprint, and task ROI. Programmatic rules derive presence, damage, count, and spatial questions with annotation- and geometry-based answers; automated checks and dual review produce 160 final tasks over 44 ROIs in three events. A shared console and headless interface then record how an agent's movement changes its observation, perception evidence, and terminal answer. This data-to-world-to-benchmark construction is the principal contribution.

In 160 dual-reviewed final questions over three generation repeats, T1\_TASK reaches 75.00\% accuracy versus 72.29\% for holding and executes fewer reobservations than the other active policies. This is an observed improvement in the tested sample, not an established stable advantage over no reobservation: the prespecified adjusted interval for the 2.71-point difference spans zero. The always-reobserve policy uses 923 actions and finishes at 63.96\%, illustrating that more observation can also be harmful. Together with the development camera-ladder's corrections and harms, these results support the platform's role in measuring task-dependent observation value.

The benchmark makes differing evidence needs explicit: damage assessment can benefit from target detail, while counts and spatial relations also depend on regional coverage and reference geometry. The observed task and event variation supports studying those requirements separately, not recommending one observation policy universally. The next scientific steps are a randomized common-budget comparison, fuller oracle decomposition, component-level timing, independently reviewed questions across more events, and a perception checkpoint evaluated without event exposure. Claims about real flight require external action and imaging validation. Until then, DisasterClaw provides an imagery-grounded platform and benchmark with explicitly bounded evidence, not a demonstrated general-purpose flight policy.
